\chapter{Installation}
\label{ch:installation}

\section{From PyPI}

With the 1.0.0 release, the Python package installs ready-built, with no compiler needed:

\begin{lstlisting}[language=bash]
pip install roadrunner-nitro
\end{lstlisting}

It is imported as \code{roadrunner\_nitro}. Until that release, build it from source as below. The C library is
always built from source.

\section{Building from source}

\subsection{Tools}

You need three tools. Open a terminal (macOS: the Terminal app; Windows: PowerShell) and check each:

\begin{enumerate}
  \item \textbf{git}: \code{git --version}. If it is missing: on macOS run \code{xcode-select --install}; on
        Windows install it from \url{https://git-scm.com}.
  \item \textbf{Python 3.10 or newer} (only for the Python package): \code{python3 --version}. If it is missing:
        \url{https://python.org}.
  \item \textbf{Rust}, which compiles the engine: \code{cargo --version}. If it is missing, install it from
        \url{https://rustup.rs}. On Windows, accept the offer to install the Visual Studio C++ Build Tools.
\end{enumerate}

\subsection{Download}

\begin{lstlisting}[language=bash]
git clone https://github.com/gshinkim/libroadrunner_nitro.git
cd libroadrunner_nitro
\end{lstlisting}

Run the commands below from inside this folder.

\subsection{The Python package}

Install into a fresh virtual environment, so nothing else on your computer changes. The first build takes one to
three minutes.

\begin{lstlisting}[language=bash]
python3 -m venv .venv
source .venv/bin/activate          # Windows: .venv\Scripts\activate
pip install --upgrade pip
pip install .
python -c "import roadrunner_nitro; print('ok')"
\end{lstlisting}

Activate the environment again (\code{source .venv/bin/activate}) in each new terminal. Fitting with pandas
tables and plotting results need \code{pandas} and \code{matplotlib}, which are optional.

\subsection{The C library}

\begin{lstlisting}[language=bash]
cargo build --release -p roadrunner_nitro_c_api
\end{lstlisting}

This makes one file in \code{target/release/}: \code{libroadrunner\_nitro\_c\_api.dylib} on macOS,
\code{libroadrunner\_nitro\_c\_api.so} on Linux and \code{roadrunner\_nitro\_c\_api.dll} (with its import
library \code{roadrunner\_nitro\_c\_api.dll.lib}) on Windows. The headers are in \code{capi/include/}.

The library is tuned for the processor of the computer that built it. To use it on another computer, delete the
\code{.cargo} folder first and build again; that makes a library that runs on any computer of the same type.
